\clearpage
\fancyhead[L]{\small FADENINTERVALLE.}
\hypertarget{native-450-fadenintervalle}{}
\begin{center}
{\large FADENINTERVALLE.}
\end{center}

\begin{center}
\rule{0.5\linewidth}{0.4pt}\\[2pt]
Astronomische Nachrichten. Band II. Nr. 43. Seite 371\,..\,376. 1823 November.\\[2pt]
\rule{0.5\linewidth}{0.4pt}
\end{center}

\begin{center}
Neue Methode, die gegenseitigen Abstände der Fäden in Meridian-Fernröhren zu bestimmen.
\end{center}

Die Bestimmungsart der Fadenintervalle in Meridian-Fernröhren, welche ausschliesslich im Gebrauch zu
sein scheint, nemlich aus beobachteten Durchgängen des Nordsterns, oder anderer dem Pole
\SourceLeaf{451}%
naher Sterne von bekannter Declination, vereinigt zwar Bequemlichkeit und grosse Schärfe, in so fern
sie unter günstigen Umständen ausgeübt werden kann. In den Jahreszeiten, wo ungünstiges Wetter
herrschend zu sein pflegt, muss man oft lange warten und manchen vergeblichen Versuch machen, bis
man nur Einen vollständigen Durchgang beobachten kann, und wenn die Gestirne, wie es so häufig der
Fall ist, nicht ganz ruhig, sondern mit starkem Zittern durch das Gesichtsfeld gehen, ist auch das
Verfahren nur einer viel geringern Genauigkeit fähig. Diese Rücksichten würden indessen viel von
ihrer Wichtigkeit verlieren, wenn man immer die definitive Bestimmung der Intervalle bis auf die Zeit
verschieben dürfte, wo die Umstände zur Anwendung jener Methode ganz vorzüglich günstig sind. Allein
theils wünscht man doch gern, die Reductionen, wozu die Kenntniss der Fadenintervalle nöthig ist, früh
machen zu können, theils kann auch das Aufschieben gefährlich werden, da besonders ein Netz von
zarten Spinnenfäden so leicht beschädigt werden kann. Ich selbst bin zuweilen in den Fall gekommen,
wegen Schlaffwerdens oder zufälligen Zerreissens eines Fadens das Netz in kurzen Fristen wiederholt
erneuern zu müssen, selbst ehe auch nur ein einziger vollständiger Durchgang des Nordsterns bei
günstiger Luft beobachtet werden konnte. Auf die Identität der Distanzen, auch wenn die neuen Fäden
in dieselben vorgerissenen Einschnitte gelegt werden, ist durchaus nicht zu rechnen, wie mich die
Erfahrung öfters gelehrt hat. Diese Gründe haben mich veranlasst, auf ein neues Mittel zur Bestimmung
der Fadenintervalle zu denken, welches neben dem gewöhnlichen, und bei ungünstigen Umständen anstatt
desselben mit Vortheil angewandt werden könnte. Dasjenige, womit ich vor kurzem eine Probe angestellt
habe, scheint dem Zwecke zu entsprechen.

So wie alle Lichtstrahlen, die unter sich parallel das Objectiv treffen, sich in Einem Punkte in der
Ebene, in welcher das Netz sich befinden soll, vereinigen, so werden umgekehrt alle Strahlen, welche
in entgegengesetzter Richtung von Einem Punkte dieser Ebene ausgehen und das Objectiv treffen, nach
dem Durchgange durch dasselbe parallel; die von verschiedenen Punkten ausgegangenen hingegen
bekommen nach dem Durchgange durch das Objectiv genau dieselben Neigungen gegen einander, die der
Entfernung jener Punkte von einander, wie sie beim Gebrauch des Instruments in der Form eines
Winkels anzusetzen ist, gleich ist. Wenn also die Ocularseite des Fernrohrs gegen den Himmel oder
sonst gegen eine helle Fläche gekehrt ist, so würde ein weitsichtiges Auge durch das Objectiv das
Fadennetz deutlich und in gehöriger Distanz (so wie eben erwähnt ist verstanden) sehen, wenn es für
so zarte Gegenstände Empfindlichkeit genug hätte: allein was dem blossen Auge unmöglich ist, wird
möglich beim Gebrauch eines zweiten Fernrohrs, wobei das Ocular eben so gestellt ist, wie es das Auge
bei sehr entfernten Gegenständen erfordert: in der That sieht man so die im Brennpunkte eines
\textsc{Frauenhofer}schen Objectivs eingezogenen Fäden mit der grössten Deutlichkeit und Reinheit. 
Auf diese Art lassen sich nun auch die Distanzen der Fäden mit grosser Präcision messen. Ich bediente
mich dazu desselben 12zolligen Theodolithen von \textsc{Ertel}, womit ich die Winkel bei der
Gradmessung beobachtet habe, und dessen eignes Netz aus einem horizontalen und zwei parallelen eine
halbe Minute von einander entfernten Fäden besteht. Bei einem neuen Ocular, welches \textsc{Repsold}
dazu verfertigt hat, sind diese Fäden von einer erstaunlichen Feinheit, und es giebt eine sehr nette
Messung, wenn man die Fäden des Meridianinstruments, durch beide Objective gesehen, mitten zwischen
jene Parallelfäden fasst.

Dass man bei dieser Operation das Ocular des Meridianfernrohrs ganz wegnehmen, oder wenigstens nur
eins mit schwächerer Vergrösserung vorsetzen muss, durch welches die Fäden, deren Intervalle gemessen
werden sollen, zugleich sichtbar sind, versteht sich von selbst.

Die Entfernung, in welcher der Theodolith von dem Objective steht, ist zwar gleichgültig; man wird
sie aber lieber nicht zu gross nehmen, um nicht einige Strahlen zu nahe am Rande des Objectivs oder
gar nur durch halbes Licht zu erhalten.

\Needspace{3\baselineskip}
\SourceLeaf{452}%
Da der Theodolith, in so fern er gehörig nivellirt ist, nicht die wirklichen, sondern die auf den
Horizont projicirten Winkel gibt, so müssen die unmittelbar gemessenen Werthe noch auf die wahren
Winkel reducirt werden. Bei so kleinen Distanzen, wie an meinen Meridianinstrumenten, ist zu diesem
Zweck die blosse Multiplication mit dem Cosinus des Neigungswinkels hinlänglich. Ist von einem
Meridiankreise die Rede, so gibt dieser denselben schon mit grösster Schärfe. Bei einem blossen
Mittagsfernrohr, welches die Neigung höchstens auf eine Minute genau angibt, kann man sie eben so
gut von dem kleinen Höhenkreise des Theodolithen entlehnen, und wird dann wohl thun, eine zu grosse
Neigung zu vermeiden.

Die folgenden Messungen sind nur als ein erster Probeversuch zu betrachten; es sind dabei noch
keinesweges alle Vorsichtsregeln angewandt, welche zu beobachten rathsam sein wird. Der Theodolith
stand mit dem hölzernen Stativ bloss auf dem bretternen Fauxplancher zwischen den Pfeilern der
Meridianinstrumente nicht unabhängig vom Stuhl des Beobachters, dessen Stellung bei der
beträchtlichen Neigung etwas genirt war, und das Licht war zum Ablesen nicht besonders günstig.
Bei Vermeidung dieser Inconvenienzen wird man daher eher eine noch bessere Uebereinstimmung zu
erwarten befugt sein. Jede Zahl ist allemal das Resultat einer 10maligen Repetition. 

\emph{Distanzen des mittelsten Fadens im Fernrohre des \textup{\textsc{Reichenbach}}schen Meridiankreises,
von den 6 übrigen.}

\begin{center}
1823 Oct.\ 23. \quad Neigung $=23^\circ\,58'\,52''$

\begingroup\renewcommand{\arraystretch}{1.1}
\begin{tabular}{r@{\;}c@{}|r|r}
 & & Horizontale & Wahre Dist. \\
\cline{3-4}
\smash{\raisebox{-0.6em}{I}} & \smash{\raisebox{-0.6em}{$\bigg\{$}} & $11'\,25''050$ & $10'\,25''93$ \\
 & & $11\quad23.000$ & $10\quad24.86$ \\
\smash{\raisebox{-0.6em}{II}} & \smash{\raisebox{-0.6em}{$\bigg\{$}} & $7\quad45.900$ & $7\quad5.69$ \\
 & & $7\quad46.275$ & $7\quad6.03$ \\
\smash{\raisebox{-0.6em}{III}} & \smash{\raisebox{-0.6em}{$\bigg\{$}} & $3\quad58.225$ & $3\quad37.66$ \\
 & & $3\quad57.750$ & $3\quad37.23$ \\
\smash{\raisebox{-0.6em}{V}} & \smash{\raisebox{-0.6em}{$\bigg\{$}} & $3\quad58.850$ & $3\quad38.24$ \\
 & & $3\quad59.375$ & $3\quad38.71$ \\
\smash{\raisebox{-0.6em}{VI}} & \smash{\raisebox{-0.6em}{$\bigg\{$}} & $7\quad53.625$ & $7\quad12.74$ \\
 & & $7\quad53.875$ & $7\quad12.98$ \\
\smash{\raisebox{-0.6em}{VII}} & \smash{\raisebox{-0.6em}{$\bigg\{$}} & $11\quad17.600$ & $10\quad19.11$ \\
 & & $11\quad17.775$ & $10\quad19.28$ \\
\end{tabular}
\endgroup

\medskip
Nov.\ 6. \quad Neigung $=23^\circ\,57'\,30''$

\begingroup\renewcommand{\arraystretch}{1.1}
\begin{tabular}{r@{\;}c@{}|r|r}
\smash{\raisebox{-0.6em}{I}} & \smash{\raisebox{-0.6em}{$\bigg\{$}} & $11'\,24''800$ & $10'\,25''29$ \\
 & & $11\quad25.575$ & $10\quad26.51$ \\
\end{tabular}
\endgroup
\end{center}

Die Fadenintervalle wurden auch aus drei beobachteten Durchgängen des Nordsterns berechnet, wovon
jedoch nur der eine vollständig hatte beobachtet werden können. Der Zustand der Luft war dabei
weder besonders günstig noch besonders ungünstig; die beiden ersten Durchgänge waren untere bei
Tage, der dritte ein oberer bei Nacht.

\begin{center}
\begingroup\small\setlength{\tabcolsep}{4pt}
\begin{tabular}{l|r|r|r|r|r|r}
 & I & II & III & V & VI & VII \\
Oct.\ 23 & $10'\,26''91$ & .\,.\,.\,. & $3'\,38''54$ & $3'\,38''12$ & $7'\,11''27$ & $10'\,19''31$ \\
Nov.\ 5 & $10\quad26.90$ & $7'\,7''01$ & $3\quad38.33$ & $3\quad39.39$ & $7\quad10.41$ & $10\quad19.50$ \\
\phantom{Nov.\ }6 & $10\quad24.31$ & .\,.\,.\,. & $3\quad35.12$ & $3\quad40.22$ & .\,.\,.\,. & .\,.\,.\,. \\
\end{tabular}
\endgroup
\end{center}

Wenn man auf die mit dem Theodolithen gemachten Bestimmungen die in meiner Theoria
Combinationis Observationum entwickelten Grundsätze anwendet, so findet sich der mittlere Fehler
eines Fadenintervalls aus 10 Repetitionen $= 0''67$; die Anwendung auf die
Nordstern-Durchgänge hingegen gibt den mittlern Fehler aus Einem Durchgange $= 1''38$, wonach
folglich die erstere Bestimmung eine doppelt so grosse Genauigkeit oder das vierfache Gewicht haben
würde. Inzwischen ist die Anzahl der
\SourceLeaf{453}%
beiderseitigen Messungen zu gering, um den mittlern Fehler mit grosser Zuverlässigkeit daraus
schliessen zu können, und in Beziehung auf den Nordstern gilt er ohnehin nur für einen ähnlichen
Luftzustand, da er bei vorzüglich schöner Luft bedeutend kleiner ausfallen würde. Im Allgemeinen aber
bestätigt sich dadurch wenigstens die Brauchbarkeit des neuen Verfahrens hinlänglich.

Eine andere nicht weniger interessante Anwendung kann man von den obigen Grundsätzen auf die
Bestimmung der Vergrösserung der Fernröhre machen, die ich jedoch hier nur kurz andeute. Insofern das
Ocular für ein weitsichtiges Auge gestellt ist, erscheinen die Gegenstände durch das Fernrohr von
hinten eben so viel verkleinert, wie dasselbe beim gewöhnlichen Gebrauche vergrössert, und zwar in
grosser Deutlichkeit, so dass sie eine Wiedervergrösserung durch ein zweites Fernrohr vollkommen
vertragen. Man kann also den Winkel zwischen zwei so betrachteten Objecten mit dem Theodolithen durch
Repetition messen, und mit dem Winkel zwischen den Richtungen von der Mitte des Oculars nach den
Objecten selbst vergleichen. Die Objecte brauchen gar nicht sehr entfernt zu sein, allenfalls, bei
etwas starker Vergrösserung, nur wenige Fuss; grösserer Schärfe wegen wird man aber, wenn die Winkel
am Ocular etwas beträchtlich sind, darauf gehörig Rücksicht zu nehmen haben, dass nicht die Winkel
selbst den Vergrösserungen proportional sind, sondern genauer, die Tangenten ihrer Neigungen gegen
die optische Axe. Man wählt also am bequemsten die Objecte so, dass zwischen den Richtungen von der
Mitte des Oculars nach denselben die Richtung der optischen Axe sehr nahe in die Mitte fällt, und
vergleicht die Tangente des halben Winkels am Ocular mit der Tangente des halben Winkels, der aus der
Beobachtung mit dem Theodolithen nach gehöriger Reduction auf den wahren Werth hervorgeht. Anstatt des
Theodolithen kann man sich auch des Heliometers bedienen, der unmittelbar den wahren Werth angibt. 
Dieses Verfahren ist einer ungleich grössern Schärfe fähig, als die andern, welche man sonst zu
gleichem Zwecke angewandt hat.
